\chapter{协作决策主干与统一通信框架}
\echapter{Collaborative Decision Backbone and Unified Communication Framework}

\section{问题定义与设计目标}
\esection{Problem Definition and Design Objectives}

本文关注的是协作型多智能体部分可观测决策问题。设系统中共有 $N$ 个智能体，在时刻 $t$ 时，智能体 $i$ 的局部观测、动作与隐藏状态分别记为 $o_i^t$、$a_i^t$ 和 $h_i^t$，联合动作为 $\mathbf{a}^t=(a_1^t,\dots,a_N^t)$，环境返回共享奖励 $r_t$。在 CTDE 范式下，训练阶段允许 Critic 使用更完整的全局信息，执行阶段则要求每个智能体仅依赖局部观测和历史记忆进行决策。

本文算法设计围绕如下三个目标展开。
\begin{enumerate}
    \item 构建一个足够稳定的 MAPPO 主干基线，使策略网络能够在无通信条件下形成可信的局部微操能力；
    \item 在不破坏主干基线的前提下，引入轻量、低风险、可解释的通信分支；
    \item 通过版本化设计逐步分析“传什么”“听谁的”和“如何融合”这三个核心问题。
\end{enumerate}

基于上述目标，本文不采用“大规模重构 Actor 主干”的方式引入通信，而是将通信建模为附着在已有策略网络之上的残差修正模块。

\section{MAPPO 协作决策主干}
\esection{MAPPO Collaborative Decision Backbone}

\subsection{局部策略网络}
\esubsection{Local Policy Network}

本文使用共享参数的循环策略网络作为 Actor。对任意智能体 $i$，在时刻 $t$ 的局部表示首先由编码器和循环单元提取：
\begin{equation}
e_i^t = f_{\text{enc}}(o_i^t), \qquad
h_i^t = \operatorname{GRU}(e_i^t, h_i^{t-1}).
\end{equation}
随后通过策略头输出局部动作 logits：
\begin{equation}
\ell_i^{\text{local}} = f_{\pi}(h_i^t),
\end{equation}
并据此定义分布式执行阶段的局部策略
\begin{equation}
\pi_{\theta}(a_i^t|o_i^t,h_i^{t-1})=
\operatorname{Softmax}\!\left(\ell_i^{\text{local}}\right).
\end{equation}

在无通信条件下，该策略网络已经能够基于局部观测学习出基础攻击、追击、后撤等微操行为，因此后续通信模块不应替代该能力，而应以小幅修正的形式提供额外协同信息。

\subsection{集中式价值函数}
\esubsection{Centralized Value Function}

Critic 使用集中式输入进行价值估计。记训练阶段可用的集中信息为 $x^t$，则集中式 Critic（centralized critic）定义为
\begin{equation}
V_{\phi}(x^t) = f_{V}(x^t).
\end{equation}
在实际实现中，$x^t$ 可由全局状态、联合观测或经过拼接的集中式特征构成。该设计的目的在于缓解部分可观测与多智能体策略耦合带来的估值困难。

基于 Critic 输出，本文采用广义优势估计（GAE）构造优势函数：
\begin{equation}
\delta_t = r_t + \gamma V_{\phi}(x^{t+1}) - V_{\phi}(x^t),
\end{equation}
\begin{equation}
\hat{A}_t = \sum_{l=0}^{T-t-1} (\gamma \lambda)^l \delta_{t+l},
\end{equation}
其中 $\gamma$ 为折扣因子，$\lambda$ 为 GAE 系数。

\subsection{MAPPO 优化目标}
\esubsection{MAPPO Objective}

本文采用 PPO 风格的截断策略目标优化共享 Actor。记新旧策略的概率比为
\begin{equation}
r_t(\theta) =
\frac{\pi_{\theta}(a_t|o_t,h_{t-1})}
{\pi_{\theta_{\text{old}}}(a_t|o_t,h_{t-1})},
\end{equation}
则策略损失为
\begin{equation}
L_{\text{policy}}(\theta)=
\mathbb{E}_t \left[
\min\left(
r_t(\theta)\hat{A}_t,
\operatorname{clip}(r_t(\theta),1-\epsilon,1+\epsilon)\hat{A}_t
\right)
\right].
\end{equation}

考虑到价值函数过快更新也会导致训练不稳定，本文进一步采用 PPO-style value clipping：
\begin{equation}
V_t^{\text{clip}} =
V_{\phi_{\text{old}}}(x_t) +
\operatorname{clip}
\left(
V_{\phi}(x_t)-V_{\phi_{\text{old}}}(x_t),
-\epsilon_v,\epsilon_v
\right),
\end{equation}
\begin{equation}
L_{\text{value}}(\phi)=
\mathbb{E}_t \left[
\max\left(
(V_{\phi}(x_t)-R_t)^2,
(V_t^{\text{clip}}-R_t)^2
\right)
\right],
\end{equation}
其中 $R_t$ 表示回报目标。最终优化目标写为
\begin{equation}
L =
L_{\text{policy}}
 - c_e L_{\text{entropy}}
 + c_v L_{\text{value}},
\end{equation}
其中 $L_{\text{entropy}}$ 为策略熵正则项，$c_e$ 与 $c_v$ 分别为熵系数和价值损失系数。

在实现层面，本文还引入了 active masks、proper time-limit、线性学习率衰减以及基于近似 KL 的 early stop 机制，从而显著增强了主干训练的稳定性。

\section{统一通信建模框架}
\esection{Unified Communication Modeling Framework}

\subsection{消息构造与路由计算}
\esubsection{Message Construction and Routing}

本文后续所有通信版本都建立在同一抽象框架之上。对智能体 $i$ 而言，首先从本地隐藏状态或其派生特征中构造待发送表示 $z_i^t$：
\begin{equation}
z_i^t = f_{\text{msg}}(h_i^t).
\end{equation}
随后分别构造 query、key 和 value：
\begin{equation}
q_i^t = W_q h_i^t,\qquad
k_j^t = W_k z_j^t,\qquad
v_j^t = W_v z_j^t.
\end{equation}
若智能体 $j$ 位于 $i$ 的候选通信集合 $\mathcal{N}_i^t$ 中，则其注意力权重定义为
\begin{equation}
\alpha_{ij}^t
=
\frac{\exp\left((q_i^t)^\top k_j^t / \sqrt{d}\right)}
{\sum_{u \in \mathcal{N}_i^t}
\exp\left((q_i^t)^\top k_u^t / \sqrt{d}\right)},
\qquad j \in \mathcal{N}_i^t.
\end{equation}
进一步得到聚合消息
\begin{equation}
m_i^t=\sum_{j \in \mathcal{N}_i^t}\alpha_{ij}^t v_j^t.
\end{equation}

这一标准化表达有助于统一理解不同版本之间的差异：有些版本主要修改 $z_i^t$ 的物理语义，有些版本主要修改 $\alpha_{ij}^t$ 的稀疏性，有些版本则主要修改 $m_i^t$ 进入策略的方式。

\subsection{残差式融合}
\esubsection{Residual Fusion}

为避免通信分支淹没局部观测特征，本文统一采用残差式融合思路。设 $\operatorname{LN}(\cdot)$ 表示归一化层，通信修正量与门控分别定义为
\begin{equation}
\Delta_i^t = f_{\Delta}\!\left([h_i^t,\operatorname{LN}(m_i^t)]\right),
\end{equation}
\begin{equation}
g_i^t = \sigma\!\left(f_g\!\left([h_i^t,\operatorname{LN}(m_i^t)]\right)+b_g\right),
\end{equation}
其中 $[\cdot,\cdot]$ 表示向量拼接，$\sigma(\cdot)$ 为 Sigmoid 函数，$b_g$ 为门控偏置。最终动作 logits 为
\begin{equation}
\tilde{\ell}_i^t
=
\ell_i^{\text{local}} + \lambda g_i^t \Delta_i^t,
\label{eq:general_residual_fusion}
\end{equation}
其中 $\lambda$ 为通信增益系数。

式 \eqref{eq:general_residual_fusion} 是全文通信设计的核心。它保证了当通信尚未学好时，策略仍可退化为纯局部 MAPPO；而当通信确实携带高价值协同信息时，门控和残差修正会逐步将其注入决策过程。

\subsection{通信设计原则}
\esubsection{Communication Design Principles}

基于上述统一框架，本文将通信设计归纳为以下四项原则。
\begin{enumerate}
    \item \textbf{低带宽}：消息维度应尽可能小，以抑制无约束高容量通信带来的数值失衡与过拟合。
    \item \textbf{低增益}：通信支路对策略的影响应从弱开始，通过零初始化和小增益逐步学习。
    \item \textbf{可解释}：优先传递具有明确物理意义的消息，例如动作意图或攻击偏置。
    \item \textbf{子空间融合}：通信最好只作用于真正需要协同的动作子空间，而不是替代整个局部策略。
\end{enumerate}

\subsection{执行期通信代价度量}
\esubsection{Execution-Time Communication Cost}

除了比较胜率与稳定性之外，本文还希望回答另一个问题：当前 strongest mainline 是否依赖于高带宽通信。为此，本文不直接使用网络参数量或训练期额外分支数作为通信开销，而是定义一个更贴近分布式执行阶段的 proxy，即每个 agent-step 的期望执行期通信维数
\begin{equation}
C_{\mathrm{exec}} = p_{\mathrm{comm}} \cdot k \cdot d_{\mathrm{exec}},
\label{eq:comm_cost_exec}
\end{equation}
其中 $p_{\mathrm{comm}}$ 表示通信支路在当前 step 保持非沉默的概率，$k$ 表示该 step 实际激活的接收者数，$d_{\mathrm{exec}}$ 表示单条 sender-receiver 边在执行期真正送入决策支路的消息维度。

式 \eqref{eq:comm_cost_exec} 的核心思想是，只统计“真正进入别的智能体决策过程的通信载体”，而不把 query、key、内部 latent 参数或训练期辅助损失都混入同一口径。这样可以统一比较三类常见情况：其一，dense all-to-all hidden-state broadcast，此时 $p_{\mathrm{comm}}=1$ 且 $k=N-1$；其二，带 top-$k$ 选择的稀疏通信，此时 $k \ll N-1$；其三，具有显式 no-comm token 或 hard gate 的 on-demand communication，此时 $p_{\mathrm{comm}}<1$。

对于本文的攻移双流通信，进一步采用逐流求和的总预算定义。若攻击流和移动流的执行期通信代价分别为
\begin{equation}
C_{\mathrm{attack}} = p_{\mathrm{attack}}\cdot k_{\mathrm{attack}}\cdot d_{\mathrm{attack}},
\end{equation}
\begin{equation}
C_{\mathrm{move}} = p_{\mathrm{move}}\cdot k_{\mathrm{move}}\cdot d_{\mathrm{move}},
\end{equation}
则总执行期通信预算写为
\begin{equation}
C_{\mathrm{total}} = C_{\mathrm{attack}} + C_{\mathrm{move}}.
\label{eq:comm_cost_total}
\end{equation}
这一写法的好处在于，它既允许分别分析 attack stream 与 move stream 的机制角色，也能在与单通道方法比较时给出一个完整的 overall communication budget。

为了让不同方法更容易横向比较，本文进一步定义两种归一化代价。第一种是同消息维度下的拓扑归一化：
\begin{equation}
\rho_{\mathrm{topo}}
=
\frac{C_{\mathrm{exec}}}{(N-1)\cdot d_{\mathrm{exec}}}
=
\frac{p_{\mathrm{comm}}\cdot k}{N-1},
\end{equation}
它反映的是在给定单边消息维度时，方法是否通过选择性路由或沉默机制真正减少了激活边数。第二种是相对 $64$ 维稠密隐藏态广播的工程归一化：
\begin{equation}
\rho_{64}
=
\frac{C_{\mathrm{exec}}}{(N-1)\cdot 64},
\end{equation}
它用于回答“当前方法的执行期通信预算，相当于向全部队友广播 $64$ 维 hidden state 的多少比例”。

需要强调的是，这一度量是执行期通信预算的 proxy，而不是总系统复杂度的完整刻画。它不包含训练期 mutual information、额外推断网络或辅助损失带来的优化开销。因此，后文若将本文方法与 MAIC 一类队友建模方法进行比较，其含义应理解为“执行期通信载体的相对轻重”，而不是整个训练系统的总成本比较。

\section{协作决策与通信流程}
\esection{Collaborative Decision and Communication Pipeline}

综合以上设计，本文的协作决策与通信总体流程可概括为如下步骤：
\begin{enumerate}
    \item 每个智能体根据局部观测更新隐藏状态，并输出本地动作 logits；
    \item 若启用通信分支，则从隐藏状态或其派生特征中构造待发送消息；
    \item 通过注意力路由对邻居消息进行选择与聚合；
    \item 将聚合消息通过门控和残差映射注入策略网络；
    \item 使用 MAPPO 的策略损失、价值损失和熵正则联合优化 Actor 与 Critic。
\end{enumerate}

这一统一框架为后续轻量通信机制设计提供了共同底座，也使得不同通信版本的差异能够被清晰地归结为“消息语义、路由结构和融合位置”的变化。

\section{本章小结}
\esection{Chapter Summary}

本章从问题定义出发，首先构建了带 centralized critic 的 vanilla MAPPO 主干，并使用 PPO-style policy clipping、value clipping、GAE 与一系列稳定化技巧提升训练鲁棒性。随后，本文提出统一通信框架，将消息设计、路由计算和残差融合纳入同一数学表达，并进一步给出执行期通信代价度量，为后续通信方法的比较提供统一口径。

这一协作决策主干与统一通信框架共同构成了全文后续方法演化与实验分析的基础。下一章将在这一底座上，进一步展开轻量通信机制设计及其面向动作层修正的具体实现。
